\centerline{\bf \Huge Турбо тренировка IV}
\centerline{\bf \Huge }

\begin{flushright}
        \scriptsize{тяжело...}
\end{flushright}

\problem{1}
Сколько 10-значных последовательностей из 0 и 1 может составить Ердем, если после каждого нуля обязательно должна стоять единичка (нулю разрешено стоять на последнем месте)?

\problem{2}
У приведенных квадратных трёхчленов $P(x)$ и $Q(x)$ есть по два (не обязательно различных) корня $p_1, p_2$ и $q_1, q_2$ соответственно. Известно, что эти числа образовывают арифметическую прогрессию (в каком-то порядке). Найдите всевозможные значения дискриминанта $P(x)$ если у $Q(x)$ он равен 7.

\problem{3}
Прямая $\ell$ касается описанной окружности остроугольного неравенобедренного треугольника $A B C$ в точке $A$. Окружность с центром $B$ и радиусом $A B$ вторично пересекает прямые $\ell$ и $A C$ в точках $P$ и $Q$. Докажите, что прямая $P Q$ проходит через ортоцентр треугольника $A B C$.

\problem{4}
\textbf{Эту задачу позорно не решил РВ на II Северо-кавказском математическом турнире.}
Эльзята выбирает случайное натуральное число $n$ и каждую секунду делает с ней следующую операцию:

1) Если число чётное, делит его на два.

2) Если число нечётное, умножает на 3 и прибавляет 1.

Докажите, что какое бы число $n$ не взяла Эльзята, она всегда в ходе своих операций получит число кратное 4.

P.s. Говорить, что гипотеза Коллатца верна нельзя)